\documentclass[10pt, letterpaper]{article}
\usepackage{xcolor}
\definecolor{bleunuit}{RGB}{25, 25, 112} % ou une autre teinte de bleu nuit

\usepackage{titlesec}
\titleformat{\section}
  {\color{bleunuit}\normalfont\Large\bfseries}
  {\thesection}{1em}{}

\usepackage[most]{tcolorbox}

% Packages:
\usepackage[
    ignoreheadfoot, % set margins without considering header and footer
    top=2 cm, % seperation between body and page edge from the top
    bottom=1.9 cm, % seperation between body and page edge from the bottom
    left=2 cm, % seperation between body and page edge from the left
    right=2 cm, % seperation between body and page edge from the right
    footskip=1.0 cm, % seperation between body and footer
    % showframe % for debugging 
]{geometry} % for adjusting page geometry
\usepackage{titlesec} % for customizing section titles
\usepackage{tabularx} % for making tables with fixed width columns
\usepackage{array} % tabularx requires this
\usepackage[dvipsnames]{xcolor} % for coloring text
\definecolor{primaryColor}{RGB}{0, 79, 144} % define primary color
\usepackage{enumitem} % for customizing lists
\usepackage{fontawesome5} % for using icons
\usepackage{amsmath} % for math
\usepackage[
    pdftitle={CV Yannick G},
    pdfauthor={Yannick Garnier},
    pdfcreator={LaTeX with RenderCV},
    colorlinks=true,
    urlcolor=primaryColor
]{hyperref} % for links, metadata and bookmarks
\usepackage[pscoord]{eso-pic} % for floating text on the page
\usepackage{calc} % for calculating lengths
\usepackage{bookmark} % for bookmarks
\usepackage{lastpage} % for getting the total number of pages
\usepackage{changepage} % for one column entries (adjustwidth environment)
\usepackage{paracol} % for two and three column entries
\usepackage{ifthen} % for conditional statements
\usepackage{needspace} % for avoiding page brake right after the section title
\usepackage{iftex} % check if engine is pdflatex, xetex or luatex

% Ensure that generate pdf is machine readable/ATS parsable:
\ifPDFTeX
    \input{glyphtounicode}
    \pdfgentounicode=1
    % \usepackage[T1]{fontenc} % this breaks sb2nov
    \usepackage[utf8]{inputenc}
    \usepackage{lmodern}
\fi



% Some settings:
\AtBeginEnvironment{adjustwidth}{\partopsep0pt} % remove space before adjustwidth environment
\pagestyle{empty} % no header or footer
\setcounter{secnumdepth}{0} % no section numbering
\setlength{\parindent}{0pt} % no indentation
\setlength{\topskip}{0pt} % no top skip
\setlength{\columnsep}{0cm} % set column seperation
\makeatletter
\let\ps@customFooterStyle\ps@plain % Copy the plain style to customFooterStyle
\patchcmd{\ps@customFooterStyle}{\thepage}{
    \color{gray}\textit{}
}{}{} % replace number by desired string
\makeatother
\pagestyle{customFooterStyle}

\titleformat{\section}{\needspace{4\baselineskip}\color{bleunuit}\bfseries\large}{}{0pt}{}[\vspace{1pt}\titlerule]


\titlespacing{\section}{
    % left space:
    -1pt
}{
    % top space:
    0.3 cm
}{
    % bottom space:
    0.2 cm
} % section title spacing

\renewcommand\labelitemi{$\circ$} % custom bullet points
\newenvironment{highlights}{
    \begin{itemize}[
        topsep=0.10 cm,
        parsep=0.10 cm,
        partopsep=0pt,
        itemsep=0pt,
        leftmargin=0.4 cm + 10pt
    ]
}{
    \end{itemize}
} % new environment for highlights

\newenvironment{highlightsforbulletentries}{
    \begin{itemize}[
        topsep=0.10 cm,
        parsep=0.10 cm,
        partopsep=0pt,
        itemsep=0pt,
        leftmargin=10pt
    ]
}{
    \end{itemize}
} % new environment for highlights for bullet entries


\newenvironment{onecolentry}{
    \begin{adjustwidth}{
        0.2 cm + 0.00001 cm
    }{
        0.2 cm + 0.00001 cm
    }
}{
    \end{adjustwidth}
} % new environment for one column entries

\newenvironment{twocolentry}[2][]{
    \onecolentry
    \def\secondColumn{#2}
    \setcolumnwidth{\fill, 4.5 cm}
    \begin{paracol}{2}
}{
    \switchcolumn \raggedleft \secondColumn
    \end{paracol}
    \endonecolentry
} % new environment for two column entries

\newenvironment{header}{
    \setlength{\topsep}{0pt}\par\kern\topsep\centering\linespread{1.5}
}{
    \par\kern\topsep
} % new environment for the header

\newcommand{\placelastupdatedtext}{% \placetextbox{<horizontal pos>}{<vertical pos>}{<stuff>}
  \AddToShipoutPictureFG*{% Add <stuff> to current page foreground
    \put(
        \LenToUnit{\paperwidth-2 cm-0.2 cm+0.05cm},
        \LenToUnit{\paperheight-1.0 cm}
    ){\vtop{{\null}\makebox[0pt][c]{
        \small\color{gray}\textit{}\hspace{\widthof{}}
    }}}%
  }%
}%

% save the original href command in a new command:
\let\hrefWithoutArrow\href

% new command for external links:
\renewcommand{\href}[2]{\hrefWithoutArrow{#1}{\ifthenelse{\equal{#2}{}}{ }{#2 }\raisebox{.15ex}{\footnotesize \faExternalLink*}}}

\vspace{-300pt}
\begin{document}
    \newcommand{\AND}{\unskip
        \cleaders\copy\ANDbox\hskip\wd\ANDbox
        \ignorespaces
    }
    \newsavebox\ANDbox
    \sbox\ANDbox{}

    \placelastupdatedtext
    \begin{header}
    \vspace*{-1.45cm}
        \textbf{\fontsize{24 pt}{24 pt}\selectfont\textcolor{bleunuit} {Yannick GARNIER}}

        \vspace{0.3 cm}

        \normalsize
        \mbox{{\color{black}\footnotesize\faMapMarker*}\hspace*{0.13cm}75000 Paris, France}%
        \kern 0.25 cm%
        \AND%
        \kern 0.25 cm%
        \mbox{\hrefWithoutArrow{mailto:yannick.garnier@gmail.com}{\color{black}{\footnotesize\faEnvelope[regular]}\hspace*{0.13cm}yannick.garnier@gmail.com}}%
        \kern 0.25 cm%
        \AND%
        \kern 0.25 cm%
        \mbox{\href{https://www.linkedin.com/}{\color{black}{\footnotesize\faLinkedin}\hspace*{0.13cm}Linkedin}}%
        \kern 0.25 cm%
        \AND%
        \kern 0.25 cm%
        \mbox{\hrefWithoutArrow{tel:+33601020304}{\color{black}{\footnotesize\faPhone*}\hspace*{0.13cm}+33 6 01 02 03 04}}%
        \kern 0.25 cm%
        \AND%
        \kern 0.25 cm%
        
        \kern 0.25 cm%
     
        \kern 0.25 cm%
        
        \kern 0.25 cm% 
    \end{header}



\vspace{-0.6cm}
\begin{tcolorbox}[
    colback=gray!5,
    colframe=bleunuit,
    coltitle=black,
    boxrule=0.4pt,
    arc=2mm,
    left=4pt,
    right=4pt,
    top=2pt,
    bottom=2pt,
    fontupper=\small,
    width=\linewidth,
    sharp corners=south
]
\justifying
Étudiant en 3\textsuperscript{ème} année à l’ENSIMAG, je recherche un stage de fin d’études (6 mois) à partir de février 2026. Curieux, analytique et à l’aise dans le travail en équipe, je souhaite contribuer à des projets innovants en informatique ou mathématiques appliquées.
\end{tcolorbox}
\vspace{-0.1cm}


\section{Formations}
\begin{twocolentry}{\textit{Depuis 2023}}
    \textbf{ENSIMAG – Grenoble INP} \\
    \textit{Filière MMIS et Master of Artificial Intelligence}
\end{twocolentry}

\vspace{0.10cm}
\begin{onecolentry}
    \begin{highlights}
        \item Formation axée sur la modélisation mathématique, l'image, la simulation et l’intelligence artificielle.
    \end{highlights}
\end{onecolentry}

\vspace{0.3cm}
\begin{twocolentry}{\textit{Avril – Août 2025}}
    \textbf{Technische Universität München (TUM)} – Allemagne \\
    \textit{Échange académique}
\end{twocolentry}

\vspace{0.10cm}
\begin{onecolentry}
    \begin{highlights}
        \item Cours en fondements mathématiques pour le Machine Learning, analyse de données et Deep Learning.
    \end{highlights}
\end{onecolentry}

\vspace{0.3cm}
\begin{twocolentry}{\textit{2021 – 2023}}
    \textbf{Cycle Préparatoire Polytechnique} – INP Grenoble \\
    \textit{Cours en sciences fondamentales et ingénierie}
\end{twocolentry}


\section{Expériences professionnelles}

\begin{twocolentry}{\textit{Grenoble, France} \\
\textit{Janvier – Avril 2025}}
    \textbf{Stagiaire assistant ingénieur} \\
    \textit{STMicroelectronics}
\end{twocolentry}

\vspace{0.10cm}
\begin{onecolentry}
    \begin{highlights}
        \item Développement en C++ d’un simulateur analogique intégré dans un simulateur digital “event-driven”.
        \item Conception et implémentation complètes à partir d’un article IEEE (2019). Projet en cours de valorisation scientifique par publication.
    \end{highlights}
\end{onecolentry}

\vspace{0.3cm}
\begin{twocolentry}{\textit{Montreux, Suisse} \\
\textit{Janvier 2025}}
    \textbf{Développeur logiciel} \\
    \textit{Clinique de gastro-entérologie}
\end{twocolentry}

\vspace{0.10cm}
\begin{onecolentry}
    \begin{highlights}
        \item Développement en C d’un logiciel automatisant l’intégration d’images dans des rapports médicaux PDF.
        \item Logiciel actuellement utilisé en production.
    \end{highlights}
\end{onecolentry}

\vspace{0.3cm}
\begin{twocolentry}{\textit{Grenoble, France} \\
\textit{Juin 2024}}
    \textbf{Technicien support logiciel} \\
    \textit{Allcom}
\end{twocolentry}

\vspace{0.10cm}
\begin{onecolentry}
    \begin{highlights}
        \item Migration vers le logiciel de facturation Dolibarr et transfert des données clients.
        \item Modification des documents via la librairie TCPDF (PHP).
    \end{highlights}
\end{onecolentry}



\section{Projets académiques}


\begin{twocolentry}{\textit{Décembre 2024}}
    \textbf{Développement d’un compilateur}
\end{twocolentry}

\vspace{0.10cm}
\begin{onecolentry}
    \begin{highlights}
        \item Création d’un compilateur pour un sous-langage de Java, en équipe de 5.
    \end{highlights}
\end{onecolentry}


\vspace{0.3cm}
\begin{twocolentry}{\textit{Mai – Juin 2024}}
    \textbf{Implémentation d’un encodeur JPEG}
\end{twocolentry}

\vspace{0.10cm}
\begin{onecolentry}
    \begin{highlights}
        \item Implémentation en C d’un encodeur JPEG selon les spécifications du standard, en équipe de 3.
    \end{highlights}
\end{onecolentry}



\section{Compétences techniques}

\begin{onecolentry}
    \textbf{Langages :} Python, C, C++, Java
\end{onecolentry}

\vspace{0.2 cm}

\begin{onecolentry}
    \textbf{Machine Learning :} PyTorch, Pandas
\end{onecolentry}

\vspace{0.2 cm}

\begin{onecolentry}
    \textbf{Bases de données :} Notions de SGBD et SQL
\end{onecolentry}

\vspace{0.2 cm}

\begin{onecolentry}
    \textbf{Outils :} Git, Linux, WordPress, TCPDF (PHP)
\end{onecolentry}

\section{Langues et centres d’intérêt}

\begin{onecolentry}
    \textbf{Langues :} Français (natif), Anglais (B2 – TOEIC), Allemand (B1)
\end{onecolentry}

\vspace{0.2 cm}

\begin{onecolentry}
    \textbf{Centres d’intérêt :} Musculation, Football, VTT, Randonnée, Machine Learning
\end{onecolentry}

% \vspace{0.2 cm}

% \begin{onecolentry}
%     \textbf{Soft skills :} Travail en équipe, autonomie, esprit d’analyse
% \end{onecolentry}

\end{document}
